\chapter{Lezione 10 --- STFT, regolatore digitale, convertitore D/A e tenuta di ordine zero}

\section{Richiamo: comportamento in frequenza della tenuta di ordine zero}

Nella parte iniziale della lezione si riprende la funzione di trasferimento della
tenuta di ordine zero, che era stata introdotta come modello del convertitore
digitale-analogico.

Se la risposta impulsiva della tenuta è un rettangolo di altezza \(1\) e durata \(T\),
allora
\[
h(t)=1(t)-1(t-T),
\]
e la sua trasformata di Laplace è
\[
F(s)=\mathcal{L}\{h(t)\}
=
\frac{1-e^{-sT}}{s}.
\]

Valutando la risposta armonica:
\[
F(j\omega)=\frac{1-e^{-j\omega T}}{j\omega}.
\]

Moltiplicando e dividendo opportunamente, si ottiene la forma standard
\[
F(j\omega)
=
e^{-j\omega T/2}\,\frac{2j\sin(\omega T/2)}{j\omega}
=
T\,e^{-j\omega T/2}\,\frac{\sin(\omega T/2)}{\omega T/2}.
\]

Questa espressione mette in evidenza due fattori distinti:
\begin{itemize}
    \item un termine di ritardo puro
    \[
    e^{-j\omega T/2},
    \]
    \item un termine reale di tipo sinc
    \[
    \frac{\sin(\omega T/2)}{\omega T/2}.
    \]
\end{itemize}

\section{Modulo della tenuta di ordine zero}

Dal risultato precedente segue
\[
|F(j\omega)|=
T\left|\frac{\sin(\omega T/2)}{\omega T/2}\right|.
\]

Quindi:
\begin{itemize}
    \item per \(\omega\to 0\), il modulo tende a \(T\);
    \item ci sono zeri quando
    \[
    \sin(\omega T/2)=0
    \quad\Longrightarrow\quad
    \omega=\frac{2k\pi}{T},\qquad k\in\mathbb{Z}\setminus\{0\};
    \]
    \item il modulo ha quindi l'andamento tipico di una sinc.
\end{itemize}

Se si considera la versione normalizzata dividendo per \(T\), si ottiene
\[
\left|\frac{F(j\omega)}{T}\right|
=
\left|\frac{\sin(\omega T/2)}{\omega T/2}\right|.
\]

Alla pulsazione di Nyquist
\[
\omega_N=\frac{\pi}{T},
\]
si ha
\[
\frac{\sin(\omega_N T/2)}{\omega_N T/2}
=
\frac{\sin(\pi/2)}{\pi/2}
=
\frac{2}{\pi}.
\]

Quindi il modulo normalizzato alla Nyquist vale
\[
20\log_{10}\!\left(\frac{2}{\pi}\right)\approx -3.92\ \mathrm{dB}.
\]

Negli appunti questa attenuazione viene approssimata qualitativamente come circa \(-3\) dB;
il valore esatto è più vicino a \(-4\) dB.

\section{Fase della tenuta di ordine zero}

La fase di \(F(j\omega)\) si può scrivere come
\[
\angle F(j\omega)
=
-\omega T/2
+
\angle\!\left(\frac{\sin(\omega T/2)}{\omega T/2}\right).
\]

Poiché il termine sinc è reale:
\begin{itemize}
    \item quando
    \[
    \frac{\sin(\omega T/2)}{\omega T/2}>0,
    \]
    il suo argomento è \(0\);
    \item quando
    \[
    \frac{\sin(\omega T/2)}{\omega T/2}<0,
    \]
    il suo argomento è \(\pi\) oppure \(-\pi\), a seconda della convenzione scelta.
\end{itemize}

Per questo la fase complessiva:
\begin{itemize}
    \item decresce linearmente come \(-\omega T/2\);
    \item subisce salti di \(\pi\) in corrispondenza degli zeri del sinc.
\end{itemize}

Se si rappresenta la fase principale, il diagramma ha il classico andamento a rampe
con riavvolgimenti periodici.

\section{Approssimazione a bassa frequenza}

Per frequenze sufficientemente piccole rispetto alla Nyquist, il termine sinc vale circa \(1\):
\[
\frac{\sin(\omega T/2)}{\omega T/2}\approx 1.
\]

Quindi
\[
F(j\omega)\approx T e^{-j\omega T/2}.
\]

Nel dominio di Laplace, questo corrisponde a
\[
F(s)\approx T e^{-sT/2}.
\]

Quindi, a bassa frequenza, la tenuta di ordine zero si comporta come:
\begin{itemize}
    \item un guadagno statico \(T\);
    \item seguito da un ritardo puro pari a \(T/2\).
\end{itemize}

Questo fatto è molto importante perché mostra che il DAC non è neutro:
introduce attenuazione alle alte frequenze e soprattutto una fase aggiuntiva
che può influenzare i margini di stabilità del sistema.

\section{Short Time Fourier Transform}

Nella prima parte della lezione viene anche introdotta la \textbf{Short Time Fourier Transform} (STFT).

L'idea è la seguente:
invece di calcolare una sola trasformata di Fourier su tutto il segnale,
si prende una finestra temporale di lunghezza finita, si calcola la trasformata
su quella finestra, poi si trasla la finestra nel tempo e si ripete il procedimento.

In questo modo si ottiene una rappresentazione bidimensionale:
\begin{itemize}
    \item asse del tempo;
    \item asse della frequenza;
    \item intensità o colore proporzionali al modulo spettrale.
\end{itemize}

\subsection{Interpretazione}

La STFT permette di vedere:
\begin{itemize}
    \item quali frequenze sono presenti;
    \item in quali intervalli temporali compaiono;
    \item come i transitori modificano il contenuto spettrale nel tempo.
\end{itemize}

Per esempio:
\begin{itemize}
    \item un transitorio veloce tende a generare contenuto ad alte frequenze;
    \item un fenomeno più lento si manifesta soprattutto alle basse frequenze.
\end{itemize}

\subsection{Compromesso tempo--frequenza}

La STFT introduce un compromesso fondamentale:
\begin{itemize}
    \item se la finestra è corta, si ha buona risoluzione temporale ma scarsa risoluzione in frequenza;
    \item se la finestra è lunga, si ha buona risoluzione in frequenza ma scarsa risoluzione temporale.
\end{itemize}

Questa è un'altra manifestazione del classico compromesso tempo--frequenza.

\section{Richiamo alle wavelet}

Negli appunti si accenna anche al fatto che esistono trasformate wavelet,
che consentono rappresentazioni tempo-frequenza con risoluzione variabile.

L'idea qualitativa è che le wavelet siano particolarmente adatte
a rappresentare fenomeni localizzati e transitori,
migliorando in molti casi il compromesso rispetto alla STFT classica.

\section{Problema di controllo digitale}

La seconda parte della lezione torna sul problema del controllo digitale.

Lo schema considerato è del tipo:
\[
y^\star \longrightarrow \text{sommatore} \longrightarrow e
\longrightarrow \text{campionatore}
\longrightarrow R(z)
\longrightarrow \text{D/A}
\longrightarrow G(s)
\longrightarrow y.
\]

Il regolatore è quindi implementato in digitale, mentre l'impianto resta continuo.

\subsection{Significato di \(G(s)\)}

Negli appunti si sottolinea che il blocco \(G(s)\) rappresenta in realtà non solo l'impianto,
ma l'intera parte continua dell'anello:
\begin{itemize}
    \item impianto fisico;
    \item attuazione;
    \item amplificazione;
    \item eventualmente anche la dinamica del sensore,
    se la si ingloba nel modello standard a retroazione unitaria.
\end{itemize}

Quindi, in pratica,
\[
G(s)=\text{impianto}+\text{attuazione}+\text{amplificazione}+\text{sensore}.
\]

\section{Tempo di calcolo del regolatore}

Se il regolatore digitale \(R(z)\) è implementato su microprocessore,
esso produce una sequenza discreta \(u(k)\) a partire dai campioni dell'errore.

In linea di principio si potrebbe pensare che, campionato l'ingresso \(e(k)\),
l'uscita \(u(k)\) venga prodotta nello stesso intervallo \([kT,(k+1)T)\).
Tuttavia il tempo di calcolo non è veramente nullo.

Negli appunti si osserva che il ritardo tra campionamento dell'ingresso
e disponibilità dell'uscita:
\begin{itemize}
    \item è inferiore a \(T\);
    \item ma non è noto con precisione;
    \item e soprattutto non è costante.
\end{itemize}

Questo è indesiderabile dal punto di vista del modello.

\subsection{Scelta pratica}

Per evitare questa incertezza, si adotta una convenzione temporale più ordinata:
\begin{itemize}
    \item il regolatore usa i campioni disponibili al passo \(k\);
    \item l'uscita effettiva viene resa coerente con il passo successivo;
    \item in questo modo ingresso e uscita del regolatore sono sempre separati da un intervallo esattamente pari a \(T\).
\end{itemize}

Questa scelta rende il modello più regolare e adatto al progetto.

\section{Due strategie di progetto del controllo digitale}

Negli appunti vengono distinte due possibili strategie.

\subsection{Prima strategia}

La più completa, ma anche più complicata, consiste nel:
\begin{itemize}
    \item discretizzare l'impianto;
    \item lavorare interamente nel dominio discreto;
    \item progettare il regolatore direttamente sul modello discreto.
\end{itemize}

\subsection{Seconda strategia}

Quella su cui il corso si concentra maggiormente è:
\begin{itemize}
    \item progettare il regolatore nel continuo, ottenendo \(R(s)\);
    \item discretizzare il solo regolatore, ad esempio con Tustin;
    \item mantenere l'impianto in forma continua.
\end{itemize}

In questo approccio serve però un buon modello del convertitore D/A,
cioè proprio della tenuta di ordine zero.

\section{Convertitore D/A come sample and hold}

Il convertitore digitale-analogico viene modellato come \textbf{sample and hold},
più precisamente \textbf{zero-order hold}.

In ingresso riceve una sequenza di campioni
\[
u(k),
\]
e in uscita produce un segnale continuo a tratti costante:
\[
u_h(t)=u(k), \qquad kT \le t < (k+1)T.
\]

Quindi il valore del campione viene mantenuto costante per tutta la durata di un periodo \(T\).

\subsection{Interpretazione}

Il segnale in uscita dal D/A non è una successione di impulsi, ma una funzione continua a tratti.
Dal punto di vista fisico, è questo segnale che entra nell'impianto \(G(s)\).

\section{Risposta impulsiva della tenuta di ordine zero}

Per determinare la funzione di trasferimento della tenuta di ordine zero,
si studia la risposta a un impulso discreto in ingresso.

Se in ingresso si applica un solo campione unitario al tempo \(k=0\),
l'uscita analogica della tenuta è un rettangolo:
\[
h(t)=
\begin{cases}
1, & 0\le t < T,\\
0, & t\ge T.
\end{cases}
\]

Equivalentemente:
\[
h(t)=1(t)-1(t-T).
\]

La trasformata di Laplace della risposta impulsiva è allora
\[
F(s)=\mathcal{L}\{h(t)\}
=
\mathcal{L}\{1(t)-1(t-T)\}
=
\frac{1}{s}-\frac{e^{-sT}}{s}
=
\frac{1-e^{-sT}}{s}.
\]

Questa è la funzione di trasferimento della tenuta di ordine zero.

\section{Risposta armonica della tenuta}

Per capire l'effetto del DAC sul comportamento complessivo del sistema,
bisogna studiare la risposta armonica della sua funzione di trasferimento.

Sostituendo \(s=j\omega\), si ottiene
\[
F(j\omega)=\frac{1-e^{-j\omega T}}{j\omega}.
\]

Come già visto, questa si può riscrivere nella forma
\[
F(j\omega)=T e^{-j\omega T/2}\frac{\sin(\omega T/2)}{\omega T/2}.
\]

\subsection{Modulo}

Il modulo vale
\[
|F(j\omega)|=T\left|\frac{\sin(\omega T/2)}{\omega T/2}\right|.
\]

Quindi:
\begin{itemize}
    \item alle basse frequenze il modulo è circa costante e pari a \(T\);
    \item all'aumentare della frequenza compare l'attenuazione di tipo sinc;
    \item il modulo si annulla ai multipli di
    \[
    \omega_c=\frac{2\pi}{T}.
    \]
\end{itemize}

\subsection{Fase}

La fase è data da
\[
\angle F(j\omega)
=
-\omega T/2
+
\angle\!\left(\frac{\sin(\omega T/2)}{\omega T/2}\right).
\]

Nel primo intervallo di frequenze, dove il sinc è positivo,
la fase è approssimativamente
\[
\angle F(j\omega)\approx -\omega T/2.
\]

Quindi, a bassa frequenza, la tenuta si comporta quasi come un ritardo puro di \(T/2\).

\section{Effetto sul diagramma di Bode}

Il DAC con tenuta di ordine zero influenza il diagramma di Bode complessivo dell'impianto:
\begin{itemize}
    \item introduce attenuazione alle alte frequenze;
    \item introduce sfasamento addizionale;
    \item può quindi peggiorare i margini di fase e il comportamento dinamico.
\end{itemize}

Per questo, nel progetto di un regolatore digitale, il blocco di tenuta
non può essere trascurato se si vuole una descrizione accurata del sistema reale.

\section{Sintesi della lezione}

I punti principali della lezione sono:
\begin{enumerate}
    \item la STFT fornisce una rappresentazione tempo--frequenza del segnale;
    \item con la STFT si vede quali frequenze sono presenti e in quali intervalli temporali;
    \item esiste un compromesso inevitabile tra risoluzione temporale e risoluzione in frequenza;
    \item nel problema di controllo digitale il regolatore \(R(z)\) è implementato in forma discreta, mentre il blocco \(G(s)\) resta continuo;
    \item il segnale prodotto dal regolatore digitale deve essere convertito in analogico tramite un convertitore D/A;
    \item il D/A è modellato come tenuta di ordine zero;
    \item la sua funzione di trasferimento è
    \[
    F(s)=\frac{1-e^{-sT}}{s};
    \]
    \item la risposta armonica della tenuta è
    \[
    F(j\omega)=T e^{-j\omega T/2}\frac{\sin(\omega T/2)}{\omega T/2};
    \]
    \item a bassa frequenza il DAC si comporta approssimativamente come un guadagno \(T\) seguito da un ritardo puro \(T/2\);
    \item il blocco di tenuta introduce attenuazione e sfasamento, e quindi modifica il Bode complessivo del sistema.
\end{enumerate}